\documentclass[aspectratio=1610]{beamer}

\title{Wildcat Beamer Theme}
\date{April 2026}
\author{Aaron Wolf (Northwestern University)}

% Theme options — pattern, color, and font can be combined:
%   \usetheme[pattern=triangle, color=nudarkblue]{wildcat}
%
% pattern: facet (default), facetrand, triangle, triangleiso, trianglemesh,
%          squares, circles, hexraise, hexrand, network, quilt, plain
%
% color: any named color (nu* colors are all available after theme loads).
%        All wcprimary shades are auto-generated. Omit to keep NU purple.
% font: overleaf (TeX Gyre), installed (system fonts), or omit for default
\usetheme{wildcat}

% You change the titlegraphic to whatever you want, or comment it out to remove it.
\titlegraphic{\includegraphics[scale=0.25]{logo-northwestern.pdf}}

% % For fine-grained control over shades, override colors manually AFTER \usetheme.
% % For example, Yale Blue (#00356b) with hand-tuned shades:
% \definecolor{wcprimary}{RGB}{0,53,107}
% \definecolor{wcprimary140}{RGB}{0, 34, 70}
% \definecolor{wcprimary130}{RGB}{0, 40, 80}
% \definecolor{wcprimary120}{RGB}{0, 45, 91}
% \definecolor{wcprimary110}{RGB}{0, 50, 102}
% \definecolor{wcprimary40}{RGB}{153, 174, 196}
% \definecolor{wcprimary30}{RGB}{179, 194, 211}
% \definecolor{wcprimary20}{RGB}{204, 215, 225}
% \definecolor{wcprimary10}{RGB}{230, 235, 240}

% % Alert and example colors:
% \definecolor{wcalerted}{RGB}{189,83,25}
% \definecolor{wcexample}{RGB}{95,113,45}

% % If you want to change the general slide background color), 
% % you can use the following command:
%\setbeamercolor{background canvas}{bg=nupurple10!30}

% % Turn off section slides
% \AtBeginSection{}


\begin{document}

\begin{frame}
\titlepage
\end{frame}

% \begin{frame}{Table of Contents}
%     \tableofcontents
% \end{frame}

\begin{frame}{Introduction}
    The Wildcat theme is a Beamer theme for Northwestern University, but which can be modified easily with different colors, fonts, and even background patterns. 
    \\ ~ \\
    The theme is inspired by the \href{https://github.com/matze/mtheme}{Metropolis theme} by Matthias Vogelgesang. It incorporates the Northwestern University facet design pattern, but otherwise has a clean, simple look, and relatively few bells and whistles. It is licensed under the GNU GENERAL PUBLIC LICENSE.
\end{frame}

\begin{frame}[fragile]{Background Patterns}
    The theme ships with 12 built-in background patterns. The easiest way to switch is via a theme option:
    \scriptsize
    \begin{verbatim}
\usetheme[pattern=triangle]{wildcat}
    \end{verbatim}
    \normalsize

    Section pages automatically use the same pattern as the title page. Available pattern options:
    \begin{columns}[t]
    \column{0.5\textwidth}
    \begin{itemize}\footnotesize
        \item \texttt{facet} — NU facets (default)
        \item \texttt{facetrand} — NU facets (alternate)
        \item \texttt{triangle} — triangle mosaic
        \item \texttt{triangleiso} — isometric triangles
        \item \texttt{trianglemesh} — wireframe mesh
        \item \texttt{squares} — square grid
    \end{itemize}
    \column{0.5\textwidth}
    \begin{itemize}\footnotesize
        \item \texttt{circles} — circles
        \item \texttt{hexraise} — raised hex clusters
        \item \texttt{hexrand} — random hex shading
        \item \texttt{network} — network grid
        \item \texttt{quilt} — quilt rings
        \item \texttt{plain} — solid color
    \end{itemize}
    \end{columns}
    \vskip0.5em
    {\footnotesize See \href{https://github.com/aarondwolf/wildcat/blob/main/source/wildcat-gallery.pdf}{\texttt{wildcat-gallery.pdf}} for a visual preview of all patterns.}
\end{frame}

\begin{frame}
    \frametitle{Colors}
    The theme has a few Northwestern-specific colors defined, which you can use in your slides. These are:
    \begin{columns}
    \column{0.3\textwidth}
    \begin{itemize}
        \item[$\textcolor{nupurple}{\bullet}$] \textcolor{nupurple}{nupurple} 
        \item[$\textcolor{nupurple90}{\bullet}$] \textcolor{nupurple90}{nupurple90} 
        \item[$\textcolor{nupurple80}{\bullet}$] \textcolor{nupurple80}{nupurple80}
        \item[$\textcolor{nupurple70}{\bullet}$] \textcolor{nupurple70}{nupurple70}
        \item[$\textcolor{nupurple60}{\bullet}$] \textcolor{nupurple60}{nupurple60}
        \item[$\textcolor{nupurple50}{\bullet}$] \textcolor{nupurple50}{nupurple50}
        \item[$\textcolor{nupurple40}{\bullet}$] \textcolor{nupurple40}{nupurple40}
        \item[$\textcolor{nupurple30}{\bullet}$] \textcolor{nupurple30}{nupurple30}
        \item[$\textcolor{nupurple20}{\bullet}$] \textcolor{nupurple20}{nupurple20}
        \item[$\textcolor{nupurple10}{\bullet}$] \textcolor{nupurple10}{nupurple10}
    \end{itemize}
    \column{0.3\textwidth}
    \begin{itemize}
        \item[$\textcolor{nupurple160}{\bullet}$] \textcolor{nupurple160}{nupurple160}
        \item[$\textcolor{nupurple150}{\bullet}$] \textcolor{nupurple150}{nupurple150}
        \item[$\textcolor{nupurple140}{\bullet}$] \textcolor{nupurple140}{nupurple140}
        \item[$\textcolor{nupurple130}{\bullet}$] \textcolor{nupurple130}{nupurple130}
        \item[$\textcolor{nupurple120}{\bullet}$] \textcolor{nupurple120}{nupurple120}
        \item[$\textcolor{nupurple110}{\bullet}$] \textcolor{nupurple110}{nupurple110}
        \item[$\textcolor{nurichblack}{\bullet}$] \textcolor{nurichblack}{nurichblack}
        \item[$\textcolor{nubrightgreen}{\bullet}$] \textcolor{nubrightgreen}{nubrightgreen}
        \item[$\textcolor{nubrightteal}{\bullet}$] \textcolor{nubrightteal}{nubrightteal}
        \item[$\textcolor{nubrightblue}{\bullet}$] \textcolor{nubrightblue}{nubrightblue}
    \end{itemize}
    \column{0.3\textwidth}
    \begin{itemize}
        \item[$\textcolor{nubrightyellow}{\bullet}$] \textcolor{nubrightyellow}{nubrightyellow}
        \item[$\textcolor{nubrightorange}{\bullet}$] \textcolor{nubrightorange}{nubrightorange}
        \item[$\textcolor{nubrightred}{\bullet}$] \textcolor{nubrightred}{nubrightred}
        \item[$\textcolor{nudarkgreen}{\bullet}$] \textcolor{nudarkgreen}{nudarkgreen}
        \item[$\textcolor{nudarkteal}{\bullet}$] \textcolor{nudarkteal}{nudarkteal}
        \item[$\textcolor{nudarkblue}{\bullet}$] \textcolor{nudarkblue}{nudarkblue}
        \item[$\textcolor{nudarkyellow}{\bullet}$] \textcolor{nudarkyellow}{nudarkyellow}
        \item[$\textcolor{nudarkorange}{\bullet}$] \textcolor{nudarkorange}{nudarkorange}
        \item[$\textcolor{nudarkred}{\bullet}$] \textcolor{nudarkred}{nudarkred}
    \end{itemize}
    \end{columns}
\end{frame}

\begin{frame}[fragile]{Modifying the Primary Color}
    Pass a \texttt{color} option to \texttt{\textbackslash usetheme} — the theme auto-generates all lighter and darker shades automatically:
    \scriptsize
    \begin{verbatim}
\usetheme[color=nudarkblue]{wildcat}
    \end{verbatim}
    \normalsize
    You can combine \texttt{color} and \texttt{pattern} in a single line:
    \scriptsize
    \begin{verbatim}
\usetheme[color=nudarkblue, pattern=hexraise]{wildcat}
    \end{verbatim}
    \normalsize
    The \texttt{nu*} colors from the previous slide are all valid choices, as are the standard xcolor built-ins (\texttt{red}, \texttt{blue}, \texttt{cyan}, etc.). The alert and example colors (\texttt{wcalerted}, \texttt{wcexample}) are not affected and can be overridden separately.
\end{frame}

\begin{frame}[fragile]{Modifying the Primary Color: Custom RGB}
    The \texttt{color} option accepts any \textit{named} color — not raw RGB values. If your brand color isn't already defined, declare it first:
    \scriptsize
    \begin{verbatim}
\definecolor{myblue}{RGB}{0,53,107}
\usetheme[color=myblue]{wildcat}
    \end{verbatim}
    \normalsize
    The auto-generated shades follow a simple scale: \texttt{wcprimary10}–\texttt{40} mix the base color with white (lighter tints), and \texttt{wcprimary110}–\texttt{140} mix it with black (darker shades). These are sufficient for most use cases — if you need precise per-shade control, see the Advanced Customization section at the end of this document.
\end{frame}

\begin{frame}[fragile]
    \frametitle{Using Colors in Graphs}
    %You can use the colors in your graphs in Python, R, and Stata using the following commands:
    \begin{tblock}{Python Code}
        \footnotesize
        \begin{verbatim}
import matplotlib.pyplot as plt
import pandas as pd
plt.style.use('ggplot')
# Create colors
wcprimary = (78/255,42/255,132/255)
df = pd.read_stata("source/graphs/auto.dta")
plt.scatter(df["weight"], df["mpg"], color=wcprimary)
plt.xlabel("Weight")
plt.ylabel("MPG")
plt.savefig("source/graphs/plot-python.pdf")
plt.show()
        \end{verbatim}
    \end{tblock}
\end{frame}

\begin{frame}[fragile]
    \frametitle{Using Colors in Graphs}
    \begin{figure}
        \centering
        \includegraphics[width=0.65\textwidth]{graphs/plot-python.pdf}
        \caption{Python Graph Example}
        \label{fig:graph-python}
    \end{figure}
\end{frame}

\begin{frame}[fragile]
    \frametitle{Using Colors in Graphs}
    \begin{tblock}{R Code}
        \footnotesize
        \begin{verbatim}
library(ggplot2)
library(haven)
# Load data
df <- read_dta("source/graphs/auto.dta")
# Define custom RGB colors
wcprimary <- rgb(78/255, 42/255, 132/255)

# Plot mpg against weight, with marker color wcprimary
plot <- ggplot(df, aes(x = weight, y = mpg)) +
    geom_point(color = wcprimary) +
    labs(x = "Weight", y = "MPG")

# Save the plot to a file with a specific size
ggsave("source/graphs/plot-R.pdf", plot, width = 8, height = 5)
        \end{verbatim}
    \end{tblock}
\end{frame}



\begin{frame}[fragile]
    \frametitle{Using Colors in Graphs}
    \begin{figure}
        \centering
        \includegraphics[width=0.8\textwidth]{graphs/plot-R.pdf}
        \caption{R Graph Example}
        \label{fig:graph-R}
    \end{figure}
\end{frame}

\begin{frame}[fragile]
    \frametitle{Using Colors in Graphs}
    \begin{tblock}{Stata Code}
        \footnotesize
        \begin{verbatim}
*ssc install schemepack, replace
sysuse auto, clear
local wcprimary "78 42 132"

twoway (scatter mpg weight, mcolor("`wcprimary'")), ///
        scheme(gg_tableau)
        
graph export "plot-stata.pdf", replace
        \end{verbatim}
    \end{tblock}
\end{frame}

\begin{frame}[fragile]
    \frametitle{Using Colors in Graphs}
    \begin{figure}
        \centering
        \includegraphics[width=0.7\textwidth]{graphs/plot-stata.pdf}
        \caption{Stata Graph Example}
        \label{fig:graph-stata}
    \end{figure}
\end{frame}


\begin{frame}{Pattern Blocks (tcolorbox)}
    You can use tcolorbox style blocks with the background pattern instead of the default beamer blocks. You can create this with a \texttt{\textbackslash begin\{tblock\}} environment. You can also use \texttt{talert} and \texttt{texample} for alert and example blocks, respectively.
    \begin{tblock}{T Block Title}
        This is a tcolorbox style block.
    \end{tblock}
    \begin{talert}{T Alert Title}
        This is a tcolorbox style alert block.
    \end{talert}
    \begin{texample}{T Example Title}
        This is a tcolorbox style example block.
    \end{texample}
\end{frame}

\begin{frame}[fragile]{Custom Color Pattern Blocks}
    There is also a special block called \texttt{tpatternbox} which allows you to specify the color. This only works with non-primary (not red, green, or blue) colors, as you can't shade those easily.
    \begin{tpatternbox}[nudarkyellow]{Custom Pattern Block Title}
        This is a tpatternbox block. It can be created with the following code:
        \begin{verbatim}
\begin{tpatternbox}[nudarkyellow]{Custom Pattern Block Title}
    This is a tpatternbox block...
\end{tpatternbox}
        \end{verbatim}
    \end{tpatternbox}
\end{frame}

\begin{frame}
    \frametitle{Box Examples (Default)}
    You can also just use the Beamer default blocks in the usual way. The default is non-rounded corners, non-shaded.
    \begin{block}{Main Block}
    This is an example block
    \end{block}
    \begin{alertblock}{Alert Box}
    This is an alert box
    \end{alertblock}
    \begin{exampleblock}{Example Box}
        This is an example box
    \end{exampleblock}
    \end{frame}

\begin{frame}[fragile]{Fonts}
    The theme uses Campton (titles) and Akkurat Pro (body) by default — these are proprietary fonts that must be purchased separately. Place the OTF files in a \texttt{fonts/} folder next to your \texttt{.tex} file:
    \scriptsize
    \begin{verbatim}
your-presentation/
  your-slides.tex
  fonts/
    Campton/
      Campton Medium.otf
      Campton Light.otf
    Akkurat/
      AkkuratProLight-Regular.otf
      AkkuratProLight-Italic.otf
      AkkuratPro-Regular.otf
      AkkuratPro-Bold.otf
    \end{verbatim}
    \normalsize
    Requires \textbf{XeLaTeX} or \textbf{LuaLaTeX} — PDFLaTeX will fail.
\end{frame}

\begin{frame}[fragile]{Fonts: Overleaf / No Proprietary Fonts}
    If you are on Overleaf or do not have the proprietary fonts, use the \texttt{font=overleaf} option. This switches to TeX Gyre fonts, which are universally available with no files to upload:
    \scriptsize
    \begin{verbatim}
\usetheme[font=overleaf]{wildcat}
    \end{verbatim}
    \normalsize
    You can combine all three options in one line:
    \scriptsize
    \begin{verbatim}
\usetheme[font=overleaf, color=nudarkblue, pattern=triangle]{wildcat}
    \end{verbatim}
    \normalsize
    For advanced font customization (loading fonts manually, using system-installed fonts), see the Advanced Customization section at the end of this document.
\end{frame}

\begin{frame}{Font Style: Examples}
    Here is what the body font will look like under normal usage:
    \begin{itemize}
        \item Regular
        \item \textit{Italic}
        \item \textbf{Bold}
        \item \textbf{\textit{Bold Italic}}
        \item \alert{Alert}
        \item \alert{\textit{Alert Italic}}
        \item Math: $$e = \lim_{n \rightarrow \infty}\left(1+\frac{1}{n}\right)^{n}$$
    \end{itemize}
\end{frame}

\begin{frame}{Sections}
    You can use \texttt{\textbackslash section\{\}} to group slides into sections. This will automatically add a section title slide at the beginning of each section.
    \\ ~ \\
    You can turn this off by adding \texttt{\textbackslash AtBeginSection\{\}} to your preamble (this tells beamer to do precisely nothing when a new section starts).
\end{frame}


\section{Advanced Customization}

\begin{frame}[fragile]{Advanced: Custom Fonts}
    If you have Campton and Akkurat Pro installed system-wide, use:
    \scriptsize
    \begin{verbatim}
\usetheme[font=installed]{wildcat}
    \end{verbatim}
    \normalsize
    To load OTF files from a custom path (e.g.\ when uploading to Overleaf), add these lines to your preamble \textit{after} \texttt{\textbackslash usetheme}:
    \footnotesize
    \begin{verbatim}
\newfontfamily\CamptonMedium[Path=fonts/Campton/, Extension=.otf]{Campton Medium}
\newfontfamily\CamptonLight[Path=fonts/Campton/, Extension=.otf]{Campton Light}
\setbeamerfont{title}{family=\CamptonMedium, size=\LARGE}
\setbeamerfont{author}{family=\CamptonLight, size=\small}
\setbeamerfont{frametitle}{family=\CamptonMedium, size=\Large}
\setbeamerfont{section}{family=\CamptonLight, size=\LARGE}
\setsansfont[Path=fonts/Akkurat/, UprightFont=*Light-Regular,
    ItalicFont=*Light-Italic, BoldFont=*-Regular,
    BoldItalicFont=*-Bold, Extension=.otf]{AkkuratPro}
    \end{verbatim}
\end{frame}

\begin{frame}[fragile]{Advanced: Fine-Grained Color Control}
    The auto-generated shades are a good starting point, but for precise brand colors you can override individual shades manually. Redefine any of the \texttt{wcprimary*} colors \textit{after} \texttt{\textbackslash usetheme}:
    \scriptsize
    \begin{verbatim}
\usetheme{wildcat}                          % or with color= option
\definecolor{wcprimary}{RGB}{0,53,107}      % Yale Blue
\definecolor{wcprimary140}{RGB}{0, 34, 70}
\definecolor{wcprimary130}{RGB}{0, 40, 80}
\definecolor{wcprimary120}{RGB}{0, 45, 91}
\definecolor{wcprimary110}{RGB}{0, 50, 102}
\definecolor{wcprimary40}{RGB}{153, 174, 196}
\definecolor{wcprimary30}{RGB}{179, 194, 211}
\definecolor{wcprimary20}{RGB}{204, 215, 225}
\definecolor{wcprimary10}{RGB}{230, 235, 240}
    \end{verbatim}
    \normalsize
    Only the shades actually used by the theme need redefining: \texttt{wcprimary10}–\texttt{40} (block bodies) and \texttt{wcprimary110}–\texttt{140} (facet pattern segments).
\end{frame}

\begin{frame}[fragile]{Advanced: Custom Background Pattern (I)}
    To create your own pattern, define a command that draws TikZ commands filling the slide. The command should take one argument — the color — so that the theme can reuse it in different contexts. Here is a minimal example:
    \scriptsize
    \begin{verbatim}
\newcommand{\mypattern}[1]{
    \fill[color=#1] (0,0) rectangle (\paperwidth,\paperheight);
}
    \end{verbatim}
    \normalsize
    The pattern command is called inside a \texttt{tikzpicture} environment whose bounding box already covers the full slide, so you only need to provide the drawing commands — no \texttt{\textbackslash begin\{tikzpicture\}} wrapper is needed.
    \\ ~ \\
    Any valid TikZ commands work. The built-in patterns each take one color argument and use \texttt{\textbackslash the\textbackslash paperwidth} / \texttt{\textbackslash the\textbackslash paperheight} to scale to the slide dimensions automatically.
\end{frame}

\begin{frame}[fragile]{Advanced: Custom Background Pattern (II)}
    Once your pattern command is defined, register it with the theme using two lines in your preamble \textit{after} \texttt{\textbackslash usetheme\{wildcat\}}:
    \scriptsize
    \begin{verbatim}
\renewcommand{\bgpattern}{\mypattern{wcprimary}}
\renewcommand{\bgpatternblock}[1]{\mypattern{#1}}
    \end{verbatim}
    \normalsize
    Two commands are needed because the theme uses the pattern in two different ways:
    \begin{itemize}
        \item \texttt{\textbackslash bgpattern} — title page, section slides, and standout frames. The color is fixed, so you bake it in when registering.
        \item \texttt{\textbackslash bgpatternblock\{color\}} — pattern block headers (\texttt{tblock}, \texttt{talert}, \texttt{texample}, \texttt{tpatternbox}), where each block supplies its own color at render time.
    \end{itemize}
\end{frame}

\begin{frame}{Advanced: Custom Background Pattern (III)}
    The 12 built-in pattern definitions in \texttt{beamerinnerthemewildcat.sty} (lines 31–386) are a useful reference and starting point.
    \\ ~ \\
    \texttt{\textbackslash bgpatternplain} (line 384) is the simplest — a single \texttt{\textbackslash fill} command — and is a good template to copy and modify. The other patterns show progressively more complex TikZ techniques: coordinate scaling, \texttt{\textbackslash foreach} loops, random seeds, scopes, and drop shadows.
    \\ ~ \\
    The \texttt{\textbackslash bgpatternblock} variant used in \texttt{tblock} headers is typically identical to \texttt{\textbackslash bgpattern}, just registered separately so the theme can pass the block's color through as an argument.
\end{frame}

% Blank Frame
\begin{frame}[plain]{}
    Here is a blank frame. You can use this to add a blank page to your presentation. You can specify this via:
    \\ ~ \\
    \texttt{\textbackslash begin\{frame\}[plain]\{\}} \\
    \texttt{...} \\
    \texttt{\textbackslash end\{frame\}}
\end{frame}

% Standout frame
\begin{frame}{Standout Slides}
    You can use the \texttt{standout} command to create a slide with a large font and no title. This is useful for a slide that you want to stand out, but which doesn't need a title.
    \\ ~ \\
    \texttt{\textbackslash standout\{Questions?\}}

\end{frame}

\standout{Questions?}


\end{document}